\scanpage{11}

$X$ topological surface

$\mathcal G$ discrete group operating on $X$

$G$ invariant subgroup of $\mathcal G$, operating
properly on $X$, so that \deleted{$\forall x\in X$}
\margin{\deleted{\illegible{1}}}
\inserted{\unsure{the}} stabilizers $G_x$ $(x\in X)$
\unsure{operate} faithfully \unsure{in a neighborhood} of $x$, while
\unsure{preserving} the local orientation at $x$

$Y=X/G$ (topological surface), $f\colon X\to Y$

$X^{!}$ discrete \unsure{subset|subset} of $X$, stable under $\mathcal G$,

$d_X$ ``signature'' on $X$ with support $X^{!}$, invariant under $\mathcal G$
\[
\unsure{i}\mapsto d_X(i)
\qquad X^{!}\longmapsto\mathbb N
\]

$Y^{!}\simeq X^{!}/G\subset Y$ \quad image of $X^{!}$ in $Y$

We assume $x\in X$, $G_x\ne\{1\}\ \Rightarrow\ x\in X^{!}$

\[
d_Y\colon Y^{!}\to\mathbb N,\qquad d_Y(y)=\nu_x d_X(x)
\]
\[
\text{if } y=f(x),\quad \nu_x=[G_x]
\]

$\mathrm{Rev}(X,d_X,\mathcal G)$ category of ramified $\mathcal G$-coverings
of $X$, \unsure{subordinate} to $d_X$

$\mathrm{Rev}(Y,d_Y,\mathcal H)$ category of ramified $\mathcal H$-\unsure{coverings}
of $Y$, \unsure{subordinate} to $d_Y$

For $X'\in\mathrm{Ob}\,\mathrm{Rev}(X,d_X,\mathcal G)$, let us consider
$X'/G$ as a surface on which
$\mathcal H=\mathcal G/G$ operates, over $Y=\unsure{X'/G|X/G}$. Then

\margin{10}

\scanpage{12}
\mspage{294}

 a) $X'/G\in\mathrm{Ob}\,\mathrm{Rev}(Y,d_Y,\mathcal H)$

 b) We thus \unsure{find|give} a functor
\[
\mathrm{Rev}(X,d_X,\mathcal G)\longrightarrow
\mathrm{Rev}(Y,d_Y,\mathcal H) \tag{19}
\]

 c) This functor is an equivalence of categories.

 d) Description of a quasi-inverse functor by ``normalized'' inverse image.

Remarks. We may assume $d_X$ with values $>0$ i.e. $d_Y$ with values
$>0$ (otherwise, if $S=\{x\in X^{!}\mid d_X(x)=0\}$, \unsure{let us replace}
$X$ by $X\setminus S$), $Y$ by $Y\setminus T$ $(T=S/G)$.

\inserted{We may assume $d_X(x)=1\ \Rightarrow\ [G_x]\ne1$
(otherwise, \unsure{replace} $X^{!}$ \unsure{by} $X^{!}\setminus S$
$S=\{x\in X^{!}\mid d_X(x)=1,\ \unsure{[G_x]=1\}}$\unsure{)}}

The most interesting case \unsure{for us is}
the one where \deleted{$d_X(x)=1$} $\forall x\in X^{!}$, \unsure{we have}
$d_X(\unsure{x})=1$ and $G_x\ne\{1\}$
(thus $X^{!}=\{x\in X\mid G_x\ne\{1\}\}$), \unsure{thus|in} the case where
\[
\mathrm{Rev}(X,d_X,\mathcal G)=\mathrm{Rev}(X,\mathcal G) \tag{20}
\]
and where $\mathrm{Rev}(\unsure{\deleted{\illegible{1}}}Y,d_Y,\mathcal H)$ can therefore
\unsure{be interpreted} as a $\mathrm{Rev}(X,\mathcal G)$
(without ramification \unsure{data})

\margin{11}

\scanpage{13}

\uline{Corollary} Let $x\in X\setminus X^{!}$, $y$ its image in
$Y\setminus Y^{!}$, we have a canonical isomorphism
\[
\Pi_1\bigl((X,d_X,\mathcal G)\unsure{,|;}x\bigr)
\overset{\sim}{\longrightarrow}
\Pi_1\bigl((Y,d_Y,\mathcal H)\unsure{,|;}y\bigr) \tag{21}
\]

We are going to \unsure{make explicit} this isomorphism in terms of
\unsure{calculations} of \deleted{loops}\inserted{paths}.

An element\inserted{$\alpha$} of
$\Pi_1\bigl((X,d_X,\mathcal G)\unsure{,|;}x\bigr)$
% ``provient'': struck through, then restored by a row of dots under it (stet).
comes from a pair $(\gamma,\ell)$, with
\[
\gamma\in\mathcal G,\qquad
\ell\ \text{class of paths in }X\setminus X^{!}
\text{ from }x\text{ to }\gamma^{-1}x
\]

(\unsure{with|for} $(\gamma,\ell)$ and $(\gamma',\ell')$ \unsure{define} the
same element, \unsure{iff|sss} $\gamma'=\gamma$, and
\inserted{the element} $\ell'^{-1}\ell$ of $\Pi_1(X\setminus X^{!},x)$ is
in the invariant subgroup kernel of
\[
\Pi_1(X\setminus X^{!},x)\longrightarrow\Pi_1((X,d_X),x)
\]
-- i.e. the invariant subgroup generated by the
\deleted{\unsure{$k$|$\ell$}} $\lambda_a^{d_X(a)}$, where
for every $a\in X^{!}$, $\lambda_a$ denotes the \unsure{class} in
$\Pi_1(X\setminus X^{!},\unsure{x|a})$ of a ``loop from $x$ around $a$'').

Then let $\dot\gamma$ be the image of $\mathcal G$ in $\mathcal H$ --
\deleted{\unsure{$f_0$|$f$}} let
$f_U\colon X\setminus X^{!}\overset{\unsure{\mathrm{df}|\mathrm{def}}}{=}U
\longrightarrow Y\setminus Y^{!}\overset{\unsure{\mathrm{df}|\mathrm{def}}}{=}V
\simeq U/G$ be the \unsure{map}

\margin{12}

\scanpage{14}
\mspage{295}

which is canonical, then the image \inserted{$\beta$} of
\inserted{$\alpha\in\Pi_1(X,d_X,\mathcal G;x)$ \unsure{defined by}}
$(\gamma,\ell)$ \deleted{is}
in $\Pi_1(Y,d_Y,\mathcal H;y)$ is $(\dot\gamma,f_U(\ell))$

(\uline{N.B.} $f_U(\ell)$ is a class of \unsure{paths} in
$V=Y\setminus Y^!$ from $y=f_U(x)$ to
$\dot\gamma^{-1}(y)=f_U(\gamma^{-1}x)$).

\inserted{\uline{N.B.}} We see at once that if
$(\gamma,\ell)$ and $(\gamma',\ell')$ \unsure{define} the
same $\alpha$, then they also \unsure{define}
the same $\beta$, which comes from the fact that
\[
\Pi_1(f_U,a):\Pi_1(X\setminus X^!,x)
\longrightarrow\Pi_1(Y\setminus Y^!,y)
\]
maps a $\lambda_a$ into \inserted{a} $\lambda_b^{\nu(a)}$
($b=f(a)$, $\nu(a)=[G_a]=$ degree of \unsure{ramification}
of $f$ at $a$) hence $\lambda_a^{d_X(a)}$ into
$\lambda_b^{d_X(a)\nu(a)=d_Y(b)}$.

Let us return to the situation of (13)
\[
X=U_{0,3}^{\mathrm{an}},\quad
\mathcal G=\mathcal T_{0,3}
=\mathfrak S_3\times\mathbb Z/2,\quad
G=\mathcal T_{0,3}^{+}
=\mathfrak S_3
\]
\margin{$X^!=\{\deleted{\illegible{1}}2,-1,\frac12;
-\bar j,\unsure{-j}\}$}
\[
Y\simeq\mathbb P^{1\,\mathrm{an}}\setminus\{0\},\quad
Y^!=\{1,\infty\},
\]
\inserted{$\mathcal H=\mathbb Z/2\mathbb Z$,}
map of \unsure{passage} to the quotient from
$\widehat X=\mathbb P^{1\,\mathrm{an}}$ to
$\widehat Y=\mathbb P^{1\,\mathrm{an}}$ given by (13)
\[
f(z)=\frac{3^3}{2^2}\frac{z^2(z-1)^2}{(z^2-z+1)^3}
\]
\margin{13}

\scanpage{15}

To calculate
$\Pi_1(Y,\underline{d}_Y,\mathbb Z/2\mathbb Z)
=\widetilde\Pi_{0,3}'$
(carto-\unsure{graphic} \unsure{triangulated} unoriented group),
we took the base point $y=-\bar j$, so we are going to choose
an $x\in\unsure{U_{0,3}^{\mathrm{an}}}$ over $y$, such that
\[
f(x)=y=-\bar j
\]
--- there is \deleted{six \unsure{places|points}} \inserted{a}
choice in a \unsure{set} $f^{-1}(\{y\})$ with six
elements, of which \uline{three} in the northern hemisphere
(the other three in the southern hemisphere),
these are permuted \unsure{transitively} by $\rho$.
They lie \unsure{respectively} in the
spherical triangles $(-\bar j,0,\frac12)$, $(-\bar j,1,2)$,
$(-\bar j,\infty,-1)$ (mapped by $f$
\unsure{homeomorphically} onto the \unsure{upper} spherical triangle $(0,1,\infty)$,
\inserted{\unsure{in the upper hemisphere)}}
with center $-\bar j$).
The choice of $\unsure{x|a}$ thus \unsure{amounts}
to the choice of one of the three
\inserted{spherical triangles, or again of the 3}
spherical segments $(0,\frac12)$, $(1,2)$, $(\infty,-1)$
--- I feel like taking $(0,\frac12)$, \unsure{corresponding}
to the triangle $(-\bar j,0,\frac12)$.

The three \inserted{involutive} generators
\deleted{\illegible{1}} $\tau'_0,\tau'_1,\tau'_\infty$
of $\Pi_1(Y,d_Y,
\unsure{\mathbb Z/2\mathbb Z|\mathbb Z/2})$
\unsure{over|over} $\tau\in\mathbb Z/2\mathbb Z$
are obtained by taking paths $\ell'_0,\ell'_1,\ell'_\infty$
from \inserted{\unsure{north pole}} $y=-\bar j$ to
\inserted{\unsure{the south pole}} $\tau y=-j$, which are \unsure{the|some}
\unsure{great circle} arcs crossing the equator resp. at $2,-1,\frac12$.

\margin{14}

\scanpage{16}
\mspage{296}

These paths are the images under $f$ of paths
on $X$, let them be $\lambda_0:\overrightarrow{x,z_0}$,
$\inserted{\lambda_1:}\overrightarrow{x,z_1}$,
$\inserted{\lambda_\infty:}\overrightarrow{x,z_\infty}$,
which \inserted{each} cross the barycentric subdivision of the double
spherical triangle $(0,1,\infty)$ in exactly one point,
\unsure{situated} \unsure{respectively} on the side
$(\frac12,-\bar j)$, $(-\bar j,0)$ and $(0,\frac12)$
of the spherical triangle $(0,\frac12,-\bar j)$.

\margin{\uline{N.B.} The points $z_0,z_1,z_\infty$
are \unsure{characterized} \inserted{\unsure{resp.$^{\mathrm t}$}}
as the only points of $X$ over
$-j=f(\unsure{\dot\tau|\tau}x)$, lying \unsure{in|in} the
spherical triangles $(-\bar j,\frac12,1)$, $(0,-\bar j,-1)$,
$(\frac12,0,-j)$.}

The points \deleted{\unsure{the}} $z_0,z_1,z_\infty$ over
$\tau y$ are \unsure{respectively} of the form
$\inserted{g_0x=}\dot\sigma_\infty\dot\tau x$,
$\inserted{g_1x=}\dot\sigma_0\dot\tau x$,
$\inserted{g_\infty x=}\deleted{\unsure{\dot\sigma_\infty}}\dot\tau x$
\inserted{(i.e. $g_0=\dot\sigma_\infty\dot\tau$,
$g_1=\dot\sigma_0\dot\tau$, $g_\infty=\dot\tau$)}.
Thus the involutive generators
$\tau'_0,\tau'_1,\tau'_\infty$ of
$\widetilde\Pi_{0,3}'=\Pi_1(Y,d_Y,\mathbb Z/2\mathbb Z;y)$
are the images of the elements
\deleted{$(\dot\sigma_{\unsure{0}}\dot\tau,\lambda_0)$},
(denoted by the same sign)
\[
\tau'_0=(\dot\sigma_0\dot\tau,\lambda_0);\quad
\tau'_1=(\dot\sigma_1\dot\tau,\lambda_1),\quad
\tau'_\infty=(\dot\sigma_\infty\dot\tau,\lambda_\infty)
\]
of $\Pi_1(X,\mathcal G;x)$.
We are going to use the great circle arc \inserted{$\ell$}
from \deleted{$x=$} $-\bar j=P$
(\deleted{origin} \inserted{\unsure{base point}}
\unsure{usual|usual} \unsure{for} the calculation of $\Pi_1(X,\mathcal G)$)
\inserted{to $x$}, to identify $\Pi_1(X,\mathcal G;x)$
with $\Pi_1(X,\mathcal G;P)$.

We can then take \deleted{\unsure{as}}
\[
\tag{21}
\tau'_0=(\dot\sigma_0\dot\tau,\lambda'_0),\quad
\tau'_1=(\dot\sigma_1\dot\tau,\lambda'_1),\quad
\tau'_\infty=(\dot\sigma_\infty\dot\tau,\lambda'_\infty)
\]
where $\lambda'_0,\lambda'_1,\lambda'_\infty$ are the
\deleted{great circle arcs}
\inserted{\deleted{\unsure{loops}} paths from}
\deleted{(i.e. the \unsure{loops|ends} at $P$)}
from $P=-\bar j$
\inserted{to \deleted{\illegible{1}} $\unsure{g_iP|\dot g_iP}$
\deleted{\illegible{1}}}
\deleted{to $z_0,z_1,z_\infty$ \unsure{respectively}
(which, with the excep-}

\margin{15}